% Generated draft: requires mathematical and source review.
\documentclass[11pt]{article}
\usepackage[margin=25mm]{geometry}
\usepackage{fontspec}
\setmainfont{Latin Modern Roman}
\usepackage{amsmath,amssymb,mathtools,unicode-math}
\setmathfont{Latin Modern Math}
\usepackage{enumitem,needspace,xurl,hyperref}
\setlength{\parindent}{0pt}
\setlength{\parskip}{5pt}
\setlength{\emergencystretch}{3em}
\newcommand{\N}{\mathbb{N}}
\newcommand{\Z}{\mathbb{Z}}
\newcommand{\Q}{\mathbb{Q}}
\newcommand{\R}{\mathbb{R}}
\newcommand{\C}{\mathbb{C}}
\newcommand{\F}{\mathbb{F}}
\newcommand{\A}{\mathbb{A}}
\newcommand{\PP}{\mathbb{P}}
\newcommand{\E}{\mathbb{E}}
\newcommand{\eps}{\varepsilon}
\newcommand{\dd}{\,\mathrm{d}}
\begin{document}
\section*{UnsolvedMath ID: 30002968}
\textbf{Source problem code:} catalogue entry
\subsection{Definitions and mathematical statement}
Let \(\mathcal A\) denote the particular generalized cluster algebra of type \(C_n\) intended by the source, let \(\mathcal X(\mathcal A)\) denote its set of cluster variables, and let \(G\) denote the associated generalized snake graph referred to there. These are placeholders for the source-specific construction: the supplied record does not specify the initial seed, coefficient system, generalized exchange polynomials, admissible values of \(n\), or construction of \(G\). A subgraph \(H\subseteq G\) is connected if every pair of its vertices can be joined by a path in \(H\).

\textbf{Question as supplied.} For every \(x\in\mathcal X(\mathcal A)\), does there exist a connected subgraph \(H\subseteq G\) that represents \(x\) in the intended generalized snake-graph construction? Do analogous representations exist for every cluster variable of each affine type covered by that same construction?

\textbf{Unresolved qualifications.} The record does not define what it means for a connected subgraph to represent a cluster variable. Its background describes the target as a perfect-matching conjecture for particular generalized type-\(C\) exchange polynomials, but supplies neither those polynomials nor a matching formula, edge weights, normalization, or correspondence. A perfect matching of a graph is a set of edges incident to each vertex exactly once; this definition alone does not determine the intended cluster-variable formula. The affine types and their constructions are also unspecified. Consequently, the question above records the intended quantifiers but is not a fully specified, independently testable conjecture. It must not be interpreted as a claim for arbitrary generalized exchange polynomials, or as merely asserting the existence of some connected subgraph without a prescribed representation of \(x\).
\subsection{Short English statement}
The supplied question asks whether every cluster variable in a particular generalized type-\(C_n\) cluster algebra has a representation by a connected part of its generalized snake graph, and whether the same construction extends to the intended affine types. The exchange polynomials, graph construction, representation rule, and affine types are missing, so a precise standalone formulation requires further source information. The supplied open status and literature assessment are dated curation, not independently verified conclusions.
\subsection{Sources}
\begin{itemize}
\item Local UnsolvedMath registry, permanent ID \(30002968\), title \emph{Snake-Graph Models for Generalized Cluster Variables}, source problem code \texttt{catalogue entry}; supplied collection: \emph{Oberwolfach Reports Open Problems}. The record gives proposed year 2015 and names no proposer.
\item Supplied report citation: \emph{Mini-Workshop: Friezes (2016)}, Oberwolfach Reports, pp.\ 2859--2894, DOI \texttt{10.4171/owr/2015/48}. \url{https://doi.org/10.4171/owr/2015/48}. The citation's title year 2016 is retained alongside the 2015 proposed year and DOI year; their relationship is not resolved by the supplied record.
\item Gleitz, \emph{Generalised Cluster Algebras and Perfect Matchings}. \url{https://ems.press/content/serial-article-files/46598}.
\item Ciliberti, \emph{Cluster Expansion Formulas for Type B and C}. \url{https://arxiv.org/abs/2405.14915}.
\item Gleitz, \emph{Generalised Cluster Algebras and q-Characters}. \url{https://doi.org/10.46298/dmtcs.2479}.
\item Supplied literature assessment, checked \textbf{2026-08-21}: the record labels the problem open and describes the 2015 perfect-matching conjecture for the stated generalized type-\(C\) exchange polynomials as unresolved. It states that later perfect-matching formulas concern ordinary type \(B\) and \(C\) cluster algebras rather than this generalized conjecture. This is supplied dated curation derived from the OWR record, not an independently verified assessment; the time-sensitive status and cited claims require independent verification.
\end{itemize}
\end{document}
